\corrigesChapitre{9}{Interactions, forces et poids}
\par\noindent\begin{minipage}{\linewidth}
\begin{corrige}{1}{Contact ou distance}
\begin{questions}
\item Contact.
\item À distance : gravitation.
\end{questions}
\end{corrige}
\end{minipage}
\par\noindent\begin{minipage}{\linewidth}
\begin{corrige}{2}{Deux appareils}
\begin{questions}
\item Une balance.
\item Un dynamomètre.
\end{questions}
\end{corrige}
\end{minipage}
\par\noindent\begin{minipage}{\linewidth}
\begin{corrige}{3}{Deux unités}
\begin{questions}
\item Kilogramme, kg.
\item Newton, N.
\end{questions}
\end{corrige}
\end{minipage}
\par\noindent\begin{minipage}{\linewidth}
\begin{corrige}{4}{Direction et sens}
\begin{questions}
\item Vers le bas : sens ; vertical : direction dans le modèle local.
\item Point d’application et valeur.
\end{questions}
\pcFigureEx{figures/rev-cor-poids.pdf}{Le poids est appliqué au centre de gravité dans ce modèle, de direction verticale et de sens vers le bas.}
\end{corrige}
\end{minipage}
\par\noindent\begin{minipage}{\linewidth}
\begin{corrige}{5}{Un premier calcul}
\begin{questions}
\item $20\unite N$.
\item Non : $200\unite g=0{,}200\unite{kg}$.
\end{questions}
\end{corrige}
\end{minipage}
\par\noindent\begin{minipage}{\linewidth}
\begin{corrige}{6}{Une petite masse}
\begin{questions}
\item $P=0{,}150\times9{,}8=1{,}47\unite N$.
\item Environ $1{,}5\unite N$.
\end{questions}
\par\smallskip\noindent\textbf{Pour comprendre.}\enspace Convertir d’abord : $150\unite g=0{,}150\unite{kg}$. Ensuite $P=mg=0{,}150\times9{,}8=1{,}47\unite N$, soit environ 1,5 N.\par
\end{corrige}
\end{minipage}
\par\noindent\begin{minipage}{\linewidth}
\begin{corrige}{7}{Retrouver la masse}
\begin{questions}
\item $m=P/g=2{,}4\unite{kg}$.
\item $2\,400\unite g$.
\end{questions}
\end{corrige}
\end{minipage}
\par\noindent\begin{minipage}{\linewidth}
\begin{corrige}{8}{Changer de lieu}
\begin{questions}
\item Masse : 60 kg ; poids : 96 N.
\item 600 N sur Terre ; la masse n’a pas changé.
\end{questions}
\par\smallskip\noindent\textbf{Pour comprendre.}\enspace $P_{\rm Terre}=60\times10=600\unite N$ et $P_{\rm Lune}=60\times1{,}6=96\unite N$. Le poids est multiplié par $1{,}6/10=0{,}16$ ; la masse ne change pas.\par
\pcFigureEx{figures/rev-cor-poids-lune.pdf}{La masse reste 60 kg ; les deux poids sont comparés avec la même échelle graphique qualitative.}
\end{corrige}
\end{minipage}
\par\noindent\begin{minipage}{\linewidth}
\begin{corrige}{9}{Une échelle}
\begin{questions}
\item $12/4=3\unite{cm}$.
\item $2{,}5\times4=10\unite N$.
\end{questions}
\par\smallskip\noindent\textbf{Pour comprendre.}\enspace Une longueur se calcule en divisant la force par le nombre de newtons représentés par un centimètre : $12/4=3\unite{cm}$. Inversement, $2{,}5\times4=10\unite N$.\par
\pcFigureExact{figures/rev-cor-force-echelle.pdf}{Flèches de 3 cm et 2,5 cm à l’impression à 100 \% ; 1 cm représente 4 N.}
\end{corrige}
\end{minipage}
\par\noindent\begin{minipage}{\linewidth}
\begin{corrige}{10}{Le livre immobile}
\begin{questions}
\item La Terre et la table.
\item Poids vertical vers le bas, action du support verticale vers le haut, de même valeur.
\end{questions}
\pcFigureEx{figures/cor-livre-forces.pdf}{Ne représenter que les forces appliquées au livre ; l’action du livre sur la table concerne un autre système.}
\end{corrige}
\end{minipage}
\par\noindent\begin{minipage}{\linewidth}
\begin{corrige}{11}{Une pente mesurée}
\begin{questions}
\item $P/m=10\unite{N/kg}$ dans les deux cas.
\item Une proportionnalité entre P et m compatible avec le modèle, à la précision des mesures.
\end{questions}
\pcFigureEx{figures/rev-cor-poids-graphe.pdf}{Les deux mesures sont compatibles avec P = 10 m ; l’origine est une prévision du modèle, pas une mesure ajoutée.}
\end{corrige}
\end{minipage}
\par\noindent\begin{minipage}{\linewidth}
\begin{corrige}{12}{Une mesure négative}
\begin{questions}
\item Contrôler le zéro et l’orientation selon la notice.
\item L’erreur peut représenter une grande part de la petite valeur mesurée.
\end{questions}
\end{corrige}
\end{minipage}
\par\noindent\begin{minipage}{\linewidth}
\begin{corrige}{13}{Des forces confondues}
\begin{questions}
\item La seconde agit sur la table, pas sur le livre.
\item L’action de la table sur le livre.
\end{questions}
\pcFigureEx{figures/cor-livre-forces.pdf}{Ne représenter que les forces appliquées au livre ; l’action du livre sur la table concerne un autre système.}
\end{corrige}
\end{minipage}
\par\noindent\begin{minipage}{\linewidth}
\begin{corrige}{14}{Un virage}
\begin{questions}
\item Non, la direction de la vitesse change.
\item Oui ; une modification de mouvement ne se limite pas à accélérer ou ralentir en valeur.
\end{questions}
\pcFigureEx{figures/cor-manege.pdf}{La direction de la vitesse change dans le mouvement circulaire uniforme ; cela n’impose pas un changement de sa valeur.}
\end{corrige}
\end{minipage}
\par\noindent\begin{minipage}{\linewidth}
\begin{corrige}{15}{Une affirmation à discuter}
\begin{questions}
\item Les effets dépendent de l’ensemble des actions et du système.
\item Une masse immobile suspendue peut subir 5 N vers le bas et 5 N vers le haut.
\end{questions}
\pcFigureEx{figures/rev-cor-forces-5N.pdf}{Ces deux forces opposées agissent sur le même système ; leurs valeurs sont égales.}
\end{corrige}
\end{minipage}
\par\noindent\begin{minipage}{\linewidth}
\phantomsection\label{cor-defi-9}\pcSolutionDefi{9}{La caisse suspendue}\begin{questions}
\item $P=39{,}2\unite N$.
\item Poids vers le bas et action du câble vers le haut, même direction et valeur ; longueur 3,92 cm chacune, environ 3,9 cm au tracé.
\item Elles agissent toutes deux sur la caisse mais proviennent de deux partenaires différents : Terre et câble.
\end{questions}
\par\smallskip\noindent\textbf{Pour comprendre.}\enspace $P=4\times9{,}8=39{,}2\unite N$. À l’équilibre dans le modèle à deux forces, la force du fil a aussi cette valeur. Chaque flèche mesure $39{,}2/10=3{,}92\unite{cm}$. Les deux flèches s’appliquent à la caisse.\par
\pcFigureExact{figures/rev-cor-caisse.pdf}{Deux flèches de 3,92 cm à l’impression à 100 \% ; même système et échelle 1 cm pour 10 N.}
\end{minipage}
\par\noindent\begin{minipage}{\linewidth}
\phantomsection\label{cor-activite-9}\pcSousTitre{Pistes pour l’investigation}
Tracer P en N en fonction de m en kg. Une droite proche de l’origine et des quotients P/m voisins de 9,8 à 10 N/kg sont attendus près du sol terrestre. La masse suspendue subit le poids et l’action du dynamomètre ; si elle est immobile dans le modèle à deux forces, leurs valeurs sont égales et les sens opposés. La force sur l’appareil concerne un autre système.\par\end{minipage}
